\section{代码与模型}
我们开源了代码与模型：

\faGithub\ \textbf{\textcolor{blue}{代码：} \href{https://github.com/LCM-Lab/Elastic-Attention}{https://github.com/LCM-Lab/Elastic-Attention}}

\faDatabase\ \textbf{\textcolor{blue}{模型：} \href{https://modelscope.cn/collections/LCM\_group/Elastic-Attention}{https://modelscope.cn/collections/LCM\_group/Elastic-Attention}}

\section{相关工作}
\label{sec:related_work}

\subsection{稀疏注意力机制}
\label{subsec:train_ssa}
为缓解标准注意力的二次复杂度，现有研究大体沿两条路径发展：推理时启发式方法与训练感知稀疏化方法。
推理时启发式方法通常采用静态模式（如固定滑窗或步长）~\citep{streamingllm, TriangleMix, Beltagy2020Longformer} 来限制感受野。
为捕捉更动态的依赖关系，研究者提出了内容感知方法：token 驱逐策略依据累计重要性分数移除低信息 token~\citep{h2o, snapkv, scissorhands}；基于 kernel 的估计器则识别关键 block 以跳过冗余计算~\citep{jiang2024minference}。
此外，prefill 优化器通过重要性驱动的选择来加速长上下文处理~\citep{flexprefill, xattention, spargeattn, peng2025acceleratingprefillinglongcontextllms, ji2025salelowbitestimation}。
尽管这些启发式方法有效，但其通常依赖敏感超参数，导致跨任务鲁棒性受限。

与之相比，训练式方法将稀疏性内化到优化目标中，以对齐训练与稀疏推理。
其主要方向之一是可学习选择机制。
例如 SeerAttention~\citep{gao2024seerattention}、NSA~\citep{NSA} 与 MoBA~\citep{moba} 采用可学习门控和层级约束去逼近真实注意力模式。
为弥合稠密预训练与稀疏适配的差距，InfLLM-v2~\citep{infllmv2} 通过无参数池化引入稠密-稀疏可切换机制；DSA~\citep{deepseekai2024deepseekv32} 则使用 lightning indexer 与两阶段训练策略高效筛选 top-$k$ key-value 对。
然而，这类方法大多聚焦于固定注意力机制内部的细粒度 block 级选择，而不是根据输入复杂度动态切换注意力计算模式本身。

\subsection{混合高效架构}
\label{subsec:hybrid_arch}
为了兼顾效率与性能，混合架构通常将全注意力（FA）与线性复杂度算子进行策略性结合。
主流范式是层间混合：在标准注意力层之间交替插入线性层（如 SSM 或 RNN），以恢复联想检索能力~\citep{ku2025systemsalgorithmsconvolutionalmultihybrid,dao2024transformersssmsgeneralizedmodels}。
代表性大规模实现如 Jamba~\citep{lieber2024jambahybridtransformermambalanguage} 使用固定块比例，后续变体则通过共享全局块~\citep{glorioso2024zambacompact7bssm} 或滑动窗口~\citep{ren2024samba} 优化内存。

近期，层内混合开始出现以提升粒度控制。
例如 PruLong~\citep{bhaskar2025cache} 与 DuoAttention~\citep{duoattention} 在单层内部混合 FA 与 Streaming Sparse Attention（SSA），将不同头分配到不同模式。
LongCat~\citep{zhang2025efficient} 提出 LoZA 机制，通过用线性复杂度 SSA 替换低敏感的 MLA 模块来构建静态 ZigZag 拓扑。
但这类方法的关键限制在于：其拓扑或比例通常在推理前静态确定，缺乏对稀疏鲁棒任务与稀疏敏感任务进行动态区分的能力，导致在多样输入下资源分配可能次优。
